\section{A confinement route inspired by \texorpdfstring{\(\ell_1\)}{l1}-regularized PageRank}
\label{sec:confinement}

\subsection{What the 2026 \texorpdfstring{\(\ell_1\)}{l1} analysis establishes}

For the same smooth PageRank quadratic \(f\), define the regularized problem
\begin{equation}
    F_\rho(x)
    :=f(x)+\alpha\rho\norm{D^{1/2}x}_1.
    \label{eq:rppr-objective}
\end{equation}
Its unique minimizer \(x_\rho^\star\) is sparse, with
\begin{equation}
    \vol\bigl(\supp(x_\rho^\star)\bigr)
    \leq\frac{1}{\rho}.
    \label{eq:rppr-volume}
\end{equation}
The recent work of \citet{fountoulakis2026complexity} gives two complementary
messages.

First, standard FISTA can be asymptotically worse than ISTA in the
locality-aware work model.  On a seed-at-leaf star, the optimum remains on the
seed leaf and ISTA stays local, whereas FISTA's extrapolation activates the
high-degree center after two extrapolated steps and incurs \(\Omega(m)\)
work.  Therefore neither accelerated iteration count nor small optimal support
is sufficient to control degree-weighted work.

Second, the paper analyzes FISTA on a slightly over-regularized objective and
uses complementarity slack to separate a small core from clearly inactive
coordinates.  Under a boundary-confinement condition, its work decomposes as
\begin{equation}
    \softO\!\left(
       \frac{1}{\rho\sqrt\alpha}
       \log\frac{\alpha}{\eps_{\rm obj}}
       +\frac{\sqrt{\vol(B)}}{\rho\alpha^{3/2}}
    \right),
    \label{eq:fista-boundary-bound}
\end{equation}
where \(B\) contains the transient spurious activations.  A graph-structural
no-percolation condition ensures that iterates remain in a core plus its
vertex boundary.  These are source results for the composite RPPR problem;
they do not directly prove locality for unregularized AESP.

\subsection{An AESP no-percolation lemma}

The changing shifted systems of AESP nevertheless have a useful structural
property: outside the current support, the proximal shift vanishes.  This
allows a direct support induction under an explicit exposure condition.

Let \(\mathcal C\subseteq V\) be a proposed core containing the source,
let
\begin{equation}
    \mathcal B:=\partial\mathcal C,
    \qquad
    \mathcal H:=\mathcal C\cup\mathcal B,
    \qquad
    \operatorname{Ext}(\mathcal C):=V\setminus\mathcal H.
    \label{eq:core-boundary}
\end{equation}
Let
\begin{equation}
    d_{\min,\mathcal B}
    :=\min_{j\in\mathcal B}d_j.
    \label{eq:min-boundary-degree}
\end{equation}
For a finite burn-in horizon \(J\), suppose
\begin{equation}
    \underline\eps
    :=\min_{1\leq t\leq J}\eps_t>0,
    \qquad
    R_z
    :=\max_{\substack{1\leq t\leq J\\0\leq k\leq K_t}}
       \norm{z_t^{(k)}}_2<\infty.
    \label{eq:threshold-and-radius}
\end{equation}

\begin{assumption}[AESP no-percolation condition]
\label{ass:aesp-no-percolation}
For every \(i\in\operatorname{Ext}(\mathcal C)\),
\begin{equation}
    \frac{\abs{N(i)\cap\mathcal B}}{d_i}
    <
    \left(
      \frac{2\underline\eps}{(1-\alpha)R_z}
    \right)^2
    d_i d_{\min,\mathcal B}.
    \label{eq:aesp-no-percolation}
\end{equation}
\end{assumption}

\begin{proposition}[Confinement of AESP inner and outer iterates]
\label{prop:aesp-confinement}
Assume \cref{ass:aesp-no-percolation}, initialize
\(x^{(0)}=y^{(0)}=0\), and use thresholded LOCGD or LOCAPPR inner updates.
Then, for every \(t\leq J\) and every inner iteration \(k\),
\begin{equation}
    \supp(x^{(t)})\cup\supp(y^{(t)})
    \cup\supp(z_t^{(k)})
    \subseteq\mathcal H.
    \label{eq:aesp-support-confinement}
\end{equation}
\end{proposition}

\begin{proof}
The base point is zero and the source belongs to \(\mathcal C\).  Assume
inductively that \(y^{(t-1)}\) and \(z=z_t^{(k)}\) are supported in
\(\mathcal H\).  Fix \(i\in\operatorname{Ext}(\mathcal C)\).  Then
\(z_i=y_i^{(t-1)}=0\), \(i\neq s\), and by definition of the vertex boundary
\(i\) has no neighbor in \(\mathcal C\).  Therefore
\begin{equation}
    \grad_i h_t(z)
    =-\frac{1-\alpha}{2}
      \sum_{j\in N(i)\cap\mathcal B}
      \frac{z_j}{\sqrt{d_i d_j}}.
    \label{eq:exterior-inner-gradient}
\end{equation}
The diagonal part of \(Qz\), the source term, and the proximal term
\(\eta(z_i-y_i^{(t-1)})\) all vanish.  By Cauchy--Schwarz,
\begin{align}
  \frac{\abs{\grad_i h_t(z)}}{\sqrt{d_i}}
  &\leq
  \frac{1-\alpha}{2d_i}
  \sqrt{\frac{\abs{N(i)\cap\mathcal B}}
              {d_{\min,\mathcal B}}}
  \norm{z}_2
  \notag\\
  &\leq
  \frac{(1-\alpha)R_z}{2d_i}
  \sqrt{\frac{\abs{N(i)\cap\mathcal B}}
              {d_{\min,\mathcal B}}}
  <\underline\eps\leq\eps_t,
  \label{eq:exterior-gradient-bound}
\end{align}
where the strict inequality is exactly
\cref{eq:aesp-no-percolation}.  Thus no exterior vertex is active.  The next
inner update remains supported in \(\mathcal H\).  At the end of the stage,
\(x^{(t)}\) is supported there, and the extrapolation
\(
 y^{(t)}=x^{(t)}+\beta(x^{(t)}-x^{(t-1)})
\)
preserves the same support.  Induction over \(k\) and \(t\) completes the
proof.
\end{proof}

\subsection{From support confinement to the AESP locality factor}

Define the one-hop gradient envelope
\begin{equation}
    \Omega:=\mathcal H\cup N(\mathcal H).
    \label{eq:gradient-envelope}
\end{equation}
If \(z\) and \(y\) are supported in \(\mathcal H\), graph locality of \(Q\)
implies
\begin{equation}
    \supp(D^{1/2}\grad h_t(z))\subseteq\Omega.
    \label{eq:gradient-support-envelope}
\end{equation}

\begin{lemma}[Support-to-locality inequality for LOCGD]
\label{lem:support-to-locality}
Let \(q\) be supported in \(\Omega\), and define the thresholded active set
\begin{equation}
    S:=\left\{u:\frac{\abs{q_u}}{d_u}\geq\eps_t\right\}.
    \label{eq:thresholded-set-generic}
\end{equation}
If \(S\neq\varnothing\) and
\(
\gamma=\norm{q_S}_1/\norm{q}_1
\), then
\begin{equation}
    \frac{\vol(S)}{\gamma}
    \leq\vol(\Omega).
    \label{eq:support-to-locality}
\end{equation}
\end{lemma}

\begin{proof}
Every active density \(\abs{q_u}/d_u\) is at least \(\eps_t\), and every
inactive density inside \(\Omega\) is smaller than \(\eps_t\).  Hence the
weighted average density on \(S\) is no smaller than the weighted average over
\(\Omega\):
\[
    \frac{\norm{q_S}_1}{\vol(S)}
    \geq
    \frac{\norm{q}_1}{\vol(\Omega)}.
\]
Rearrange and use the definition of \(\gamma\).
\end{proof}

Averaging \cref{eq:support-to-locality} over LOCGD inner iterations gives
\begin{equation}
    \Lambda_t\leq\vol(\Omega).
    \label{eq:locgd-envelope-locality}
\end{equation}
For an online coordinate method, the same conclusion follows from a
normalized Gauss--Southwell rule
\begin{equation}
    u_k\in\arg\max_u\frac{\abs{q_u^{(k)}}}{d_u}.
    \label{eq:priority-rule}
\end{equation}
Indeed,
\begin{equation}
    d_{u_k}
    \leq
    \vol(\Omega)
    \frac{\abs{q_{u_k}^{(k)}}}{\norm{q^{(k)}}_1},
    \label{eq:priority-cost-charge}
\end{equation}
so the degree cost can be charged to the online active-ratio sum.  This yields
the same locality factor up to priority-queue logarithms.  Confinement alone
does not imply the corresponding statement for an arbitrary FIFO order; a
queue-order lemma would still be needed.

\begin{theorem}[Conditional end-to-end bound under confinement]
\label{thm:confinement-total}
Assume \cref{ass:aesp-no-percolation} and suppose
\begin{equation}
    \vol(\Omega)\leq\frac{K}{\eps}.
    \label{eq:small-envelope}
\end{equation}
Use AESP-LOCGD or AESP with the normalized priority rule
\cref{eq:priority-rule} until the objective-gap handoff, followed by
\(\omega=1\) LocSOR.  Then
\begin{equation}
    \Work_{\rm hybrid}
    =\softO\!\left(
       \frac{K+1}{\sqrt\alpha\eps}
    \right).
    \label{eq:confinement-total-work}
\end{equation}
\end{theorem}

\begin{proof}
By \cref{prop:aesp-confinement}, all early inner gradients are supported in
\(\Omega\).  By \cref{lem:support-to-locality} or
\cref{eq:priority-cost-charge},
\(\Lambda_t\leq\vol(\Omega)\leq K/\eps\).  Apply
\cref{thm:conditional-total}.
\end{proof}

\subsection{What this route does and does not prove}

A natural analytical choice is
\begin{equation}
    \mathcal C=\supp(x_\rho^\star),
    \qquad \rho\asymp\eps,
    \label{eq:rppr-core-choice}
\end{equation}
for which \cref{eq:rppr-volume} gives
\(\vol(\mathcal C)=\mathcal O(1/\eps)\).  This fact alone does not imply
\cref{eq:small-envelope}; a single boundary vertex can have arbitrarily high
degree.  The star lower bound of \citet{fountoulakis2026complexity} is exactly
such a warning.  Moreover, AESP in this note solves the unregularized
quadratic, whereas \(\mathcal C\) is defined through a regularized optimum.
The RPPR core is therefore a proof device, not an automatically invariant
support.

The most promising refinement is to combine AESP with slight
\emph{over-regularization}, as in the 2026 FISTA analysis.  Low-margin
coordinates would be absorbed into a larger \(\mathcal O(1/\eps)\)-volume
core, clearly inactive coordinates would be charged through KKT slack, and the
remaining boundary work would be summed using accelerated outer contraction.
Adapting that argument to AESP's sequence of changing shifted subproblems is a
substantial open proof obligation, not a completed theorem.
